%% Abstract section.
\abstract{Uncertainty visualization has become an increasingly important topic given the ubiquity of noise in data and computational processes. This workshop proposal is motivated by the uncertainty visualization challenges identified during the {\em 2024 IEEE Workshop on Uncertainty Visualization}. Specifically, four research gaps were identified. First, uncertainty must be treated as data in its own right; however, visualization community currently lacks understanding of good ways to model uncertainty and propagate it through complex scientific workflows. Second, the use of artificial intelligence (AI) is becoming pervasive across scientific domains; however, it faces a major limitation of trustworthiness of predictions due to uncertainties propagated from the data and AI models. Third, there are no standard ways to effectively communicate uncertainty with end users who can be scientific or general audience. Finally, performing decision-making under uncertainty is an open challenge, considering the role of perception, cognition, and user expertise when interpreting uncertainty. To tackle the aforementioned challenges, we propose to conduct another workshop on uncertainty visualization at IEEE VIS 2025. In particular, the overarching goal of the workshop is to create a roadmap to tackle the identified uncertainty visualization challenges that are multidisciplinary in nature by fostering interactions among experts from visualization, applications, AI, high-performance computing, applied math, and psychology fields.} % end of abstract

%The overarching goal of the workshop is to bring together this multi-disciplinary group to enlighten the visualization community in the following four areas:
%(1) How uncertainty can be efficiently represented and propagated (2) How generalized frameworks for uncertainty visualization can be developed for broader impact (3) How uncertainty in AI models can be discovered to enable reliable science, (3) How uncertainty can be effectively communicated, (4) How users should perform decision-making under uncertainty.%


%\abstract{Uncertainty visualization has become an increasingly important topic given the ubiquity of noise in data and computational processes. Although the research in uncertainty visualization has steadily progressed over the past few years, this critical branch of visualization is still in its infancy given many difficult challenges (e.g., computation, rendering, perception and decision-making) relevant to communication and understanding of uncertainty. One important step to address these challenges is to provide a venue that attracts a wide range of experts across many disciplines.
%A venue that allows experts in visualization, applications, applied math, perception, and cognition to publish and discuss effective ways to convey and understand uncertainty is an important step in advancing this critical area of research.
%The goal of the workshop is to bring together this multi-disciplinary group to enlighten the visualization community in the following four areas:
%(1) use cases in diverse application domains that can benefit from visualization of uncertainty (2) theory, techniques, and state-of-the-art software for uncertainty visualization (3) Methods/workflows that enable robust decisions under uncertainty (4) development of a future roadmap of uncertainty visualization research goals.%
%} % end of abstract

%
%not only visualization experts but also attracts experts from the  applications, applied math, and perceptual and cognitive sciences. 
%attracts application scientists, psychologists, and mathematicians who can publish and discuss effective ways to convey and understand uncertainty. Our goal of the workshop is therefore to bring together application scientists, mathematicians, and psychologists who may not be visualization experts, but can enlighten the visualization community in the following three areas: